\documentclass[12pt]{article}
\usepackage[utf8]{inputenc}
\usepackage[T1]{fontenc}
\usepackage{hyperref}
\title{Warlmanpa Language Resources for K-12 Curriculum Development}
\author{Compiled by Claude Sonnet 5.5 (claude-sonnet-5-5) for Andr\'as Kornai}
\date{}
\begin{document}
\maketitle

\section{Introduction}
Warlmanpa (ISO 639-3 wrl; Glottocode warl1255; Pama-Nyungan) was traditionally spoken northwest of Tennant Creek, Northern Territory. Most remaining speakers now live in Tennant Creek. The documentation is thin and consists mostly of a new reference grammar and older word lists.

\section{Primary Educational Resources}
I found no primers, readers or school materials specific to Warlmanpa. AUSTLANG records no literacy material for it. Mitchell Browne's grammar notes that its description draws on work by community members and linguists since 1952, but it is a reference work and not classroom material.

\section{Linguistic and Reference Materials}
\begin{itemize}
\item Mitchell Browne, \emph{A Grammar of Warlmanpa} (ANU Press, November 2024), the first published reference grammar: \url{https://press-prod.anu.edu.au/publications/series/asia-pacific-linguistics/grammar-warlmanpa}.
\item The Macquarie University research record for the grammar is at \url{https://researchers.mq.edu.au/en/publications/a-grammar-of-warlmanpa/}.
\item The AIATSIS AUSTLANG entry lists Nash 1979 as a word list and grammar sketch in a practical orthography: \url{https://collection.aiatsis.gov.au/austlang/language/C17}.
\item David Nash's Warlmanpa reference list includes a 1981 \emph{Preliminary Vocabulary of the Warlmanpa Language}. I did not confirm where a copy is held: \url{https://www0.anu.edu.au/linguistics/nash/aust/wpa/lx-ref.html}.
\item A Territory Stories item (a 1997 body-parts word list, ``Palka'', by Bessie Campbell, from the Living Archive of Aboriginal Languages): \url{https://territorystories.nt.gov.au/10070/882513}.
\end{itemize}

\section{Religious and Community Resources}
I found no Warlmanpa scripture or dedicated community organisation. Papulu Apparr-kari Aboriginal Corporation in Tennant Creek is the main regional language centre named in results for neighbouring Warumungu. It may know of Warlmanpa materials, but the results do not confirm that, and I give no contact details.

\section{Digital and Audio Resources}
I found no audio or digital resources specific to Warlmanpa.

\section{Curriculum Development Context}
Warlmanpa is not known to be taught in schools. The new grammar and the word lists are the only usable base. Curriculum work would need new materials made with speakers and the Tennant Creek regional language centre. Little exists beyond what is listed here.
\end{document}